A lo largo de este \textit{checkpoint} se complet\'o la integraci\'on, puesta a punto y validaci\'on experimental del convertidor \textit{buck} a lazo cerrado. El trabajo permiti\'o pasar de un dise\~no te\'orico y simulado a una implementaci\'on f\'isica funcional, verificando en cada etapa los bloques principales de la fuente: la etapa de potencia, el circuito de control, el \textit{gate driver}, la red de realimentaci\'on y el PCB final.

En primer lugar, la caracterizaci\'on experimental del inductor confirm\'o que el componente seleccionado era adecuado para la aplicaci\'on. La inductancia medida, de aproximadamente $408{,}7\,\upmu$H, result\'o superior al valor m\'inimo requerido por el dise\~no, mientras que la corriente de saturaci\'on de $4\,$A otorg\'o un margen suficiente frente a las corrientes esperadas durante la operaci\'on normal. Esto permiti\'o asegurar el funcionamiento del convertidor en modo de conducci\'on continua dentro del rango de carga ensayado.

La etapa de control tambi\'en fue ajustada de manera satisfactoria. Las modificaciones realizadas sobre el generador de PWM y el dise\~no de la red de compensaci\'on permitieron estabilizar el sistema a lazo cerrado y mejorar la respuesta frente a perturbaciones. Las simulaciones evidenciaron una reducci\'on de las oscilaciones respecto del sistema sin compensar y una convergencia adecuada de la tensi\'on de salida hacia el valor de referencia, con un margen de fase compatible con una operaci\'on robusta.

En la implementaci\'on pr\'actica se verific\'o el correcto funcionamiento del \textit{gate driver} y se midi\'o un tiempo muerto acorde con el especificado por el fabricante, aspecto fundamental para evitar la conducci\'on simult\'anea de los MOSFET. Adem\'as, la incorporaci\'on de diodos Schottky, capacitores apropiados para la frecuencia de conmutaci\'on y una carga fantasma mejor\'o el comportamiento del circuito durante los transitorios y facilit\'o la realizaci\'on de los ensayos experimentales.

Las mediciones sobre el prototipo confirmaron que la fuente cumple con los objetivos principales del proyecto. Para $V_{in}=12\,$V se obtuvieron eficiencias superiores al $90\%$ dentro del rango de corrientes evaluado, mientras que al variar la tensi\'on de entrada con corriente de carga constante la eficiencia se mantuvo por encima del $85\%$. Asimismo, se observ\'o una buena regulaci\'on de l\'inea y de carga: la tensi\'on de salida present\'o variaciones reducidas, del orden de $10\,$mV al duplicar aproximadamente la corriente de salida.

El desarrollo y fabricaci\'on del PCB permitieron consolidar el dise\~no en una versi\'on compacta y reproducible. Esta etapa puso en evidencia la importancia del trazado de las mallas de corriente, la ubicaci\'on de los capacitores de desacople y salida, y la selecci\'on tecnol\'ogica de los componentes. En particular, el reemplazo del capacitor de salida por uno cer\'amico permiti\'o corregir el comportamiento del \textit{ripple} y acercar la respuesta medida a la prevista te\'oricamente.

En conclusi\'on, se logr\'o dise\~nar, simular, construir y validar experimentalmente una fuente \textit{buck} a lazo cerrado. Si bien durante el proceso aparecieron diferencias entre el modelo ideal y el comportamiento real del circuito, estas fueron identificadas y corregidas mediante ajustes de dise\~no y selecci\'on de componentes. El resultado final presenta una regulaci\'on adecuada, alta eficiencia y estabilidad operativa, por lo que cumple con los requerimientos planteados y constituye una implementaci\'on funcional del sistema propuesto.